% Archival definitions unused by the three Grade 11 reference documents.
% Not loaded by the active preambles. Requires the packages/configuration
% from preamble.tex; definitions below are preserved verbatim.

% Boolean switch for optional explanations.
% Use \mostrarExplicaciones{true} to show explanations.
% Use \mostrarExplicaciones{false} to hide explanations.
\newboolean{explicaciones}

% \explicacion{shown text}{hidden-mode text}
% Displays the first argument when explanations are enabled;
% otherwise displays the second argument.
% Usage example: \explicacion{Detailed explanation.}{}
\newcommand{\explicacion}[2]{%
	\ifthenelse{\boolean{explicaciones}}{#1}{#2}%
}

% Sets the explanation switch.
% Usage example: \mostrarExplicaciones{true}
\newcommand{\mostrarExplicaciones}[1]{%
	\setboolean{explicaciones}{#1}%
}

\definecolor{popUp}{HTML}{666666}
\definecolor{popUpIn}{HTML}{CED9E0}
\definecolor{backgroundC}{HTML}{D0E8F2}
\definecolor{backgroundCC}{HTML}{FFFFFF}
\definecolor{borders}{HTML}{8c120d}
\definecolor{padding}{HTML}{77D0D7}
% General mdframed style for highlighted content blocks.
\mdfdefinestyle{estiloGeneral}{
	linecolor=black,
	linewidth=1.5pt,
	roundcorner=10pt,
	backgroundcolor=backgroundC,
	innerbottommargin=0pt
}

% mdframed style for code-like blocks.
\mdfdefinestyle{code}{
	linecolor=black,
	linewidth=1.5pt,
	roundcorner=10pt,
	backgroundcolor=darkgray!10,
	innerbottommargin=0pt
}

% Inline tcolorbox used for framed links and linked images.
\newtcbox{\fboxC}{
	colback=backgroundC,
	colframe=popUp,
	arc=10pt,
	boxrule=3pt,
	boxsep=0pt,
	nobeforeafter
}

% \cuadro{content}
% Wraps content in the standard general frame.
% Usage example: \cuadro{Important information.}
\newcommand{\cuadro}[1]{%
	\begin{mdframed}[style=estiloGeneral]
		#1
	\end{mdframed}%
}

% Python listings style.
\lstdefinestyle{python}{
	language=Python,
	basicstyle=\ttfamily\small,
	commentstyle=\color{green!50!black},
	keywordstyle=\color{blue},
	numberstyle=\tiny\color{gray},
	numbers=left,
	morekeywords={>, <},
	breakatwhitespace=false,
	showstringspaces=false,
	showtabs=false,
	showspaces=false
}

% \bigLink{url}{text}
% Creates a large centered framed hyperlink.
% Usage example: \bigLink{https://example.com}{Open resource}
\newcommand{\bigLink}[2]{%
	\begin{center}
		\fboxC{\LARGE{\href{#1}{#2}}}
	\end{center}%
}

% \smallLink{url}{text}
% Creates a small centered framed hyperlink.
% Usage example: \smallLink{https://example.com}{Open}
\newcommand{\smallLink}[2]{%
	\begin{center}
		\fboxC{\href{#1}{#2}}
	\end{center}%
}

% \bfLink{url}{text}
% Creates a bold inline hyperlink.
% Usage example: \bfLink{https://example.com}{Important link}
\newcommand{\bfLink}[2]{%
	\href{#1}{\textbf{#2}}%
}

% \smallUrl{url}
% Displays a centered framed URL.
% Usage example: \smallUrl{https://example.com}
\newcommand{\smallUrl}[1]{%
	\begin{center}
		\fboxC{\url{#1}}
	\end{center}%
}

% \subSecLink{url}{title}
% Creates an unnumbered subsubsection title as a hyperlink.
% Usage example: \subSecLink{https://example.com}{Additional practice}
\newcommand{\subSecLink}[2]{%
	\subsubsection*{\href{#1}{\textbf{#2}}}%
}

% ==========================================================
% CUSTOM LAYOUT AND PAGE COMMANDS
% ==========================================================
% \background
% Draws a large colored background rectangle behind the page content.
% Usage example: \background
\newcommand{\background}[0]{%
	\begin{tikzpicture}[remember picture,overlay]
		\fill[backgroundC] (-2,2) rectangle (25cm,-550);
	\end{tikzpicture}%
}

% \backgroundC
% Draws a large white background rectangle behind the page content.
% Kept with its original name for compatibility.
% Usage example: \backgroundC
\newcommand{\backgroundC}[0]{%
	\begin{tikzpicture}[remember picture,overlay]
		\fill[backgroundCC] (-2,2) rectangle (25cm,-550);
	\end{tikzpicture}%
}

% \anchoPag
% Stores the intended page/content width used by custom layouts.
% Usage example: \parbox{\anchoPag}{Content}
\newcommand{\anchoPag}[0]{20cm}

% General vertical spacing helper.
% Usage example: Text before.\esp Text after.
\newcommand{\esp}[0]{%
	\vspace{4mm}%
}

% ==========================================================
% CUSTOM FIGURE COMMANDS
% ==========================================================
% Horizontal offset used by image commands.
% Usage example: \hspace*{\espacioImagenes}
\newcommand{\espacioImagenes}[0]{-1.2cm}

% \fig[horizontal offset]{width fraction}{image path}
% Inserts an image without a frame.
% Usage example: \fig{0.8}{example-image}
% Usage example: \fig[-1cm]{0.8}{example-image}
\newcommand{\fig}[3][\espacioImagenes]{%
	\hspace*{#1}
	\centering
	\includegraphics[width=#2\textwidth]{#3}%
}

% \ffig{width fraction}{image path}
% Inserts a framed image inside a figure environment.
% Usage example: \ffig{0.8}{example-image}
\newcommand{\ffig}[2]{%
	\begin{figure}[h]
		\hspace*{\espacioImagenes}
		\centering
		\fbox{\includegraphics[width=#1\textwidth]{#2}}
	\end{figure}%
}

% \hfig{url}{width fraction}{image path}
% Inserts a framed linked image inside a figure environment.
% Usage example: \hfig{https://example.com}{0.8}{example-image}
\newcommand{\hfig}[3]{%
	\begin{figure}[h]
		\hspace*{-1.4cm}
		\centering
		\color{popUp}
		\fboxC{\href{#1}{\includegraphics[width=#2\textwidth]{#3}}}
	\end{figure}%
}

% \hffig{url}{width fraction}{image path}
% Inserts a linked image without a frame inside a figure environment.
% Usage example: \hffig{https://example.com}{0.8}{example-image}
\newcommand{\hffig}[3]{%
	\begin{figure}[h]
		\hspace*{-1.1cm}
		\centering
		\color{popUp}
		\href{#1}{\includegraphics[width=#2\textwidth]{#3}}
	\end{figure}%
}

% \linkExplicacion{url}
% Inserts the standard explanation-video image as a clickable figure.
% Usage example: \linkExplicacion{https://example.com/video}
\newcommand{\linkExplicacion}[1]{%
	\hffig{#1}{0.5}{principal/videoExplicacion}
	\vspace{-0.5cm}%
}



